\documentclass{article}
\usepackage[T1]{fontenc}
\usepackage{lmodern}
\usepackage{graphicx}
\usepackage{amsmath,amssymb}
\usepackage[a4paper,margin=2cm]{geometry}
\usepackage[absolute,overlay]{textpos}
\setlength{\TPHorizModule}{1mm}
\setlength{\TPVertModule}{1mm}
\usepackage[indonesian]{babel}
\usepackage{hyperref}
\setlength{\emergencystretch}{2em}
% unit: Halaman 31

\begin{document}

% === Halaman 31 (llm) ===

\textbf{Pengurangan Dua Titik di dalam EC pada GF(p)}\ \\
Misalkan $P(x_p, y_p)$ dan $Q(x_q, y_q)$. \\
Pengurangan: $P - Q = P + (-Q)$, yang dalam hal ini $-Q(x_q, -y_q \pmod{p})$.

\end{document}
